% ============================================================
% 1. INTRODUCTION
% ============================================================
\section{Introduction}
\label{sec:intro}

% --- Module 1: Why this problem matters (~45%) ---

Large language models (LLMs) have rapidly advanced across a wide
range of tasks, from question answering and code generation to web
navigation and embodied decision making~\citep{yao2023react,
shinn2023reflexion, zhou2024lats}. Yet many tasks remain beyond the
reach of a single forward pass: complex reasoning, multi-step
planning, and open-ended exploration often require the model to
evaluate alternatives, verify intermediate results, or search over
candidate actions~\citep{yao2023tree, hao2023reasoning,
snell2024scaling}. This has motivated \emph{test-time compute}, where additional
computation at inference time improves decision quality. Across
multiple benchmarks, test-time optimization allows smaller models to
match or exceed the performance of larger ones at lower total
cost~\citep{snell2024scaling, wang2026flare}.

However, test-time compute is expensive, incurring
5--15$\times$ overhead per step~\citep{yao2023tree, zhou2024lats}. As
agents are deployed across diverse settings, applying this overhead
uniformly becomes wasteful when only a fraction of steps benefit.
This has motivated \emph{selective triggering}: using uncertainty
signals to decide \emph{when} extra compute is
worthwhile~\citep{catts2026, seag2025, cats2025, corefine2026,
adaptthink2025, thinkless2025, s1_2025, han2025tokenbudget}. Across diverse mechanisms, including entropy thresholding,
calibrated confidence, vote disagreement, and reinforcement
learning, selective triggering consistently reduces cost while
preserving or improving task performance within each method's
evaluation setting.

Yet methods calibrated for one task type frequently fail when
transferred to others~\citep{tao2025revisiting,
heo2025llmuncertainty}. Vote-based
gating designed for web agents degrades on fact-verification tasks;
calibration-based methods that perform well on QA degrade in code
generation. In several cases, selective triggering \emph{actively
harms} performance compared to never invoking the optimizer. This
raises a central question: \emph{do uncertainty signals have a
universal relationship with compute value, or does this relationship
depend on the environment?}

% --- Module 2: Our finding + diagnosis (~30%) ---

We conduct systematic measurement across eight agent environments
spanning QA, code generation, web navigation, fact verification,
and text games, using three model backbones. The signal--utility
direction \emph{reverses} across environments.
The same entropy value that predicts rollout benefit in code
generation predicts rollout harm in fact verification. The same
signal in the same environment can also reverse across model
backbones. Beyond direction, even the identity of the most
informative signal varies: structural features dominate in some
environments while entropy carries no information.

We trace this reversal to a conflation of two uncertainty sources.
\emph{Information poverty} arises when the agent lacks sufficient
data; here rollouts from the same model cannot compensate, and high
entropy predicts rollout harm. \emph{Decision difficulty} arises
when the agent faces multiple viable paths; here rollouts explore
alternatives productively, and high entropy predicts rollout benefit.
Environments differ in the proportion of these two state types,
causing the same signal to carry opposite meaning. This is an
instance of Simpson's paradox in the signal--utility space. We prove
that this has a sharp consequence: no fixed-direction gate can
guarantee non-negative value of computation across environments with
opposing directions (Proposition~\ref{prop:necessity}). Once the
direction is wrong, more precise calibration makes performance
\emph{worse}, because the gate more precisely targets harmful states.

% --- Module 3: Method + contributions (~25%) ---

Motivated by this analysis, we propose \textbf{EAAG}
(\textbf{E}nvironment-\textbf{A}ware \textbf{A}daptive
\textbf{G}ating), which discovers environment-specific gating
patterns through randomized \emph{exploration}, LLM-based
\emph{reasoning} about exploration data, and sparse
\emph{learning} of signal direction and importance via LASSO.
Since direction is the primary determinant of gate quality, not
model capacity, EAAG pairs a rich discovery process with a
deliberately simple gate that adds zero per-step deployment cost.
EAAG outperforms all fixed-direction baselines across eight
environments and three LLM backbones.

Our main contributions are as follows:
\begin{itemize}[leftmargin=*,nosep]
\item \textbf{Empirical finding.} Signal--utility direction reverses
  across environments and model backbones, and the identity of the
  most informative signal varies entirely. All fixed-direction methods
  fail in at least two of eight environments.

\item \textbf{Theoretical explanation.} A two-source uncertainty
  model grounded in Simpson's paradox explains when and why direction
  reverses. Direction discovery is provably necessary for
  cross-environment non-negative value of computation.

\item \textbf{Adaptive gating method.} EAAG combines exploration,
  LLM-based reasoning, and sparse direction learning to discover
  environment-specific gating patterns with zero per-step deployment
  cost.
\end{itemize}
